\clearpage\thispagestyle{empty}
\pdfbookmark[0]{Sources for the July 1802 Pallas Report}{pallas180207-sources}
\begin{center}\textbf{Sources for the July 1802 Pallas Report}\end{center}
\noindent This AI translation follows the German excerpt in Carl Friedrich Gauss’s \href{https://gdz.sub.uni-goettingen.de/id/PPN236007467}{\textit{Werke}, vol.~VI} (Göttingen, 1874), pp.~217--220. Franz Xaver von Zach reports on Pallas, discovered by Heinrich Wilhelm Matthias Olbers. The excerpt contains calculations by Gauss and quotes his letter of 26 June.
\par\smallskip\noindent Von Zach’s \href{https://api.digitale-sammlungen.de/iiif/image/v2/bsb10538598_00077/full/1450,/0/default.jpg}{original article}, no.~X, appeared in \textit{Monatliche Correspondenz zur Beförderung der Erd- und Himmels-Kunde} (Gotha, Becker), vol.~VI, July 1802, pp.~71--96. The collected excerpt begins with the \href{https://api.digitale-sammlungen.de/iiif/image/v2/bsb10538598_00086/full/1450,/0/default.jpg}{letter on p.~80}, resumes with the third set of orbital elements on p.~83 and concludes with the \href{https://api.digitale-sammlungen.de/iiif/image/v2/bsb10538598_00100/full/1450,/0/default.jpg}{discussion of diameters on pp.~94--95}.
\par\smallskip\noindent The \href{https://api.digitale-sammlungen.de/iiif/image/v2/bsb10538598_00090/full/1450,/0/default.jpg}{journal, p.~84}, cites Ceres elements VII for the mean motion and gives \(+0{,}7\) and \(0\) as the right-ascension differences for 5 and 7 April. The collected edition cites IV and gives \(+0.\) and \(0.7\). Both translations follow the collected edition.
\par\smallskip\noindent The paper by \href{https://api.digitale-sammlungen.de/iiif/image/v2/bsb10538597_00562/full/1450,/0/default.jpg}{Prof.\ Wurm} cited in the letter is on planetary masses (June 1802, pp.~546--570). The separate paper by \href{https://doi.org/10.1098/rstl.1802.0010}{William Herschel} is \textit{Observations on the Two Lately Discovered Celestial Bodies} (1802, pp.~213--232).
\par\smallskip\noindent Translation and typesetting produced by AI. Journal scans from the Bayerische Staatsbibliothek.
\clearpage
